% Standalone fragment: requires amsmath, amssymb, amsthm, hyperref and proposition/corollary environments.
% Suggested citations are specified with verified metadata in discovery.md.
% This fragment has no dependencies on the main manuscript's theorem labels.

\paragraph{Discovery requires admission opportunities.}
The following is a standard conditional-hazard argument, closely related to
conditional Borel--Cantelli
(\href{https://sites.math.duke.edu/~rtd/PTE/PTE5_011119.pdf}{Durrett, 2019,
Theorem~4.3.4}), applied to insertion of a verified key into memory.
Its hazard includes scheduling, validation, capacity, and candidate selection;
a positive probability of generating a key alone is insufficient.

\begin{proposition}[Conditional admission hazards]
Fix a key $b$ not initially banked. Let $(\mathcal F_n)_{n\ge0}$ describe the
complete training history after each opportunity, and let $\tau_b$ be its first
admission time, with $\tau_b=\infty$ if it is never admitted. Define the
predictable stopped hazard
\[
 h_{b,n}=\Pr(\tau_b=n\mid\mathcal F_{n-1}),\qquad
 H_{b,N}=\sum_{n=1}^N h_{b,n}.
\]
For every $a\ge0$,
\[
 \Pr(\tau_b>N,\ H_{b,N}\ge a)\le e^{-a}.
\]
In particular,
$\Pr(\tau_b=\infty,\ H_{b,\infty}=\infty)=0$.
If deterministic numbers $\underline h_{b,n}\in[0,1]$ satisfy
$h_{b,n}\ge\underline h_{b,n}$ on $\{\tau_b\ge n\}$, then
\[
 \Pr(\tau_b>N)\le\prod_{n=1}^N(1-\underline h_{b,n})
 \le\exp\!\left(-\sum_{n=1}^N\underline h_{b,n}\right).
\]
For a finite target set $S$, the probability that some target key has not been
admitted by $N$ is at most the sum of these bounds over $b\in S$.
\end{proposition}
\begin{proof}
Let $Z_N=\mathbf1\{\tau_b>N\}\exp(H_{b,N})$, with $Z_0=1$.
The event of an earlier admission is known at time $n-1$, and $h_{b,n}=0$
on that event. On survival through $n-1$,
\[
 \mathbb E[Z_n\mid\mathcal F_{n-1}]
 =Z_{n-1}e^{h_{b,n}}(1-h_{b,n})\le Z_{n-1}.
\]
Thus $Z_n$ is a nonnegative supermartingale. Markov's inequality gives the
finite-hazard bound. Fatou's lemma excludes a positive-probability event on
which survival continues forever while $H_{b,N}\to\infty$, since then
$Z_N\to\infty$. Under deterministic hazard floors, conditioning gives
$\Pr(\tau_b>n)\le(1-\underline h_{b,n})\Pr(\tau_b>n-1)$;
iteration and the union bound finish the proof.
\end{proof}

Divergent cumulative hazard is a substantive premise. Independent
opportunities with admission probabilities $h_n=(n+1)^{-2}>0$ satisfy
\[
 \Pr(\tau=\infty)=\prod_{n=1}^\infty
 \left(1-\frac1{(n+1)^2}\right)=\frac12.
\]
Likewise, once a non-evicting capacity-$K$ bank is full, every unbanked key has
zero admission hazard. These bounds establish coverage only for keys that
continue to have admissible opportunities; they do not establish such
opportunities for the implemented policy or for supports larger than capacity.

\paragraph{Retention across bank changes.}
An admitted key also needs persistent loss weight. The next elementary energy
accounting result allows bank changes while explicitly charging their effect
on the objective. It is not an assumption that different bank objectives
share a single unchanged potential. This follows the standard switched-energy
viewpoint (\href{https://people.ece.ubc.ca/moishi/eece571m/articles/Branicky98.pdf}{Branicky, 1998});
the specific inequality below is proved directly.

\begin{proposition}[Switching-energy retention]
Let $t_0<t_1<\cdots$ be switching times with finitely many switches on every
bounded time interval. On interval $r$, use a finite collection of verified
complete-response exemplars, with nonnegative coefficients $d_{rj}$ and
\[
 F_r(\theta)=\sum_{j:d_{rj}>0}d_{rj}\ell_j(\theta)-J(\theta),
 \qquad \ell_j(\theta)=-\log\pi_\theta(e_j\mid x_j),
 \qquad J(\theta)\le J_{\max}<\infty.
\]
Zero coefficients denote absent exemplars and are omitted from the sum.
Suppose initial energy is finite, parameters are continuous
at switches, $F_r$ is nonincreasing within its interval, and the positive
energy jumps obey
\[
 A:=\sum_{r\ge1}
 \left[F_r(\theta(t_r))-F_{r-1}(\theta(t_r))\right]_+<\infty.
\]
If exemplar $j$ is protected from interval $r_j$ onward and
$d_{rj}\ge\underline d_j>0$ for every $r\ge r_j$, then, with
$D=F_0(\theta(t_0))+A+J_{\max}<\infty$,
\[
 p_\theta(b_j\mid x_j)\ge\pi_\theta(e_j\mid x_j)
 \ge\exp(-D/\underline d_j)
 \qquad(t\ge t_{r_j}).
\]
The likelihood must represent the complete verified response under the law
whose key probability is bounded. Mean-token losses have
$d_{rj}=\alpha_{rj}/L_j$ for fixed exemplar length $L_j>0$.
\end{proposition}
\begin{proof}
Induction over the switches gives
$F_r(\theta(t))\le F_0(\theta(t_0))+A$. Therefore
$\sum_jd_{rj}\ell_j(\theta(t))\le D$.
Each summand is nonnegative, so a protected exemplar has
$\ell_j(\theta(t))\le D/\underline d_j$.
Its complete response event is contained in its verified-key event.
\end{proof}

If old coefficients are preserved when new exemplars are added, a switch
costs exactly the summed weighted surprisal of the new exemplars at admission.
For uniform categorical replay with fixed dose $\rho$, adding one key to a
size-$k$ bank instead changes the potential by
\[
 \Delta F=\frac{\rho}{k+1}
 \bigl[-\log p_{\rm new}-R_{\rm old}\bigr].
\]
Old coefficients decrease from $\rho/k$ to $\rho/(k+1)$, but remain at least
$\rho/K$ when capacity is $K$. Finitely many such admissions, with positive
admission-time probabilities and no eviction, therefore have finite total
switching cost in the exact descent model. Recurrent scheduling, bounded
priority multipliers, or a finite-capacity buffer alone do not establish
within-interval descent or a summable jump budget for arbitrary optimization.


\begin{corollary}[Conditional discovery followed by retention]
Suppose a prompt has a finite correct support $\mathcal C$ of size $m$ no
larger than bank capacity. Admissions are non-evicting, and every missing key
$b$ has conditional admission hazard at least $\epsilon r_b$ at every
opportunity, for constants $\epsilon>0$ and $r_b>0$ with
$\epsilon r_b\le1$. Assume also the switching-energy conditions above,
including a finite positive-jump budget and persistent positive coefficients
for the admitted complete exemplars. Then the bank reaches full correct
coverage in finitely many opportunities almost surely, and thereafter every correct key has
a uniform-in-time positive probability floor. At horizon $N$,
\[
 \Pr(\text{full coverage by }N)
 \ge 1-\sum_{b\in\mathcal C}e^{-N\epsilon r_b}.
\]
If $r_b\ge\gamma>0$ for every correct key, this is at least
$1-m e^{-N\epsilon\gamma}$.
\end{corollary}
\begin{proof}
Apply the admission-hazard bound to each initially missing key and take a
union bound. The failure bound tends to zero; finite support and non-eviction
therefore imply a finite almost-sure opportunity index of simultaneous full coverage.
The switching-energy bound protects each admitted complete exemplar and hence
its key, and applies to all keys after that time.
\end{proof}

A fixed full-support proposal mixture supplies such a hazard floor if it is
queried at every counted opportunity and insertion is guaranteed whenever the
key appears. This is a sufficient design condition, not a property established
for the current temperature schedule, bounded search attempts, or adaptive
priority controller. The corollary guarantees retention, not uniform key
probabilities or convergence of the neural optimizer.
